\section{Gradient là một tích ma trận}
\label{sec:ch03-tich-jacobian}

\index{BPTT!dạng tổng tích}%
Chương~\ref{ch:ky-hieu-va-bptt} dẫn ra
\tn{lan truyền ngược qua thời gian}{backpropagation through time}, viết tắt là
BPTT, theo cách quen thuộc: đẩy ngược
từng bước một, cộng dồn vào các \tn{ma trận trọng số}{weight matrix}. Bài báo
viết lại cùng công
thức đó dưới dạng tổng tích, và chính cách viết lại này làm vấn đề lộ ra. Chi
phí toàn chuỗi tách theo từng bước,

\begin{equation}
  \pd{L}{\theta} = \sum_{1 \le t \le T} \pd{L_t}{\theta},
  \label{eq:ch03-tach-chi-phi}
\end{equation}

và phần của một bước lại tách tiếp theo mọi bước đứng trước nó:

\begin{equation}
  \pd{L_t}{\theta}
    = \sum_{1 \le k \le t}
      \pd{L_t}{\vect{a}_t}\,
      \jac{\vect{a}_t}{\vect{a}_k}\,
      \frac{\partial^{+} \vect{a}_k}{\partial \theta}.
  \label{eq:ch03-tach-theo-buoc}
\end{equation}

Dấu cộng trên \(\partial^{+}\) là ký hiệu của bài: đạo hàm riêng \emph{tức
thời} của \(\vect{a}_k\) theo \(\theta\), coi \(\vect{a}_{k-1}\) là hằng. Nó
tách phần \enquote{\(\theta\) tác động ngay tại bước \(k\)} khỏi phần
\enquote{tác động đó đi được bao xa}. Phần thứ hai là thừa số ở giữa, và nó là
toàn bộ nội dung của chương này.

Thừa số ấy là một tích:

\begin{equation}
  \jac{\vect{a}_t}{\vect{a}_k}
    = \prod_{t \ge i > k} \jac{\vect{a}_i}{\vect{a}_{i-1}}
    = \prod_{t \ge i > k} \mat{W}_{hh} \diag\!\left( \sigma'(\vect{a}_{i-1}) \right).
  \label{eq:ch03-tich-jacobian}
\end{equation}

\index{Jacobian!tích qua thời gian}%
Một tích của \(t - k\) ma trận. Bài báo có một câu về nó mà tôi thấy là câu
đáng nhớ nhất trong cả bài: giống như tích của \(t-k\) số thực có thể co về
không hoặc nổ ra vô cùng, tích ma trận này cũng vậy, nhưng \emph{theo một hướng
nào đó}. Mệnh đề trong ngoặc là chỗ trực giác vô hướng đánh rơi, và
mục~\ref{sec:ch03-khe-ho} dựa hoàn toàn vào chỗ đó.

Một lưu ý về quy ước trước khi đi tiếp. Bài in thừa số trong~\eqref{eq:ch03-tich-jacobian}
là \(\mat{W}_{hh}\transpose \diag(\sigma'(\vect{a}_{i-1}))\), vì họ đẩy gradient
ngược dưới dạng vector hàng. Chuyển vị đúng của Jacobian cột phải là
\(\diag(\sigma'(\vect{a}_{i-1})) \mat{W}_{hh}\transpose\), nên hai biểu thức
không phải chuyển vị của nhau theo nghĩa đen. Không có hệ quả nào: mọi chặn
trong bài đặt trên \tn{chuẩn phổ}{spectral norm}, và một ma trận với chuyển vị của nó có cùng chuẩn
phổ. Sách khai báo quy ước thay vì thừa kế chỗ mập mờ; phụ
lục~\ref{app:ky-hieu} ghi lại.

\subsection*{Điều kiện đủ để gradient tắt}

\index{vanishing gradient!điều kiện đủ}%
Chặn từng thừa số một. Với \(\gamma\) là chặn trên của \(\abs{\sigma'(x)}\),

\begin{equation}
  \norm{\jac{\vect{a}_i}{\vect{a}_{i-1}}}
    \le \norm{\mat{W}_{hh}} \cdot \norm{\diag(\sigma'(\vect{a}_{i-1}))}
    \le \norm{\mat{W}_{hh}} \, \gamma.
  \label{eq:ch03-chan-tung-thua-so}
\end{equation}

Nếu \(\norm{\mat{W}_{hh}} < 1/\gamma\) thì vế phải nhỏ hơn 1. Đặt \(\eta\) là
một số nào đó nằm giữa, tức \(\norm{\jac{\vect{a}_i}{\vect{a}_{i-1}}} \le \eta < 1\)
với mọi \(i\). Chuẩn của một tích không vượt quá tích các chuẩn, nên quy nạp
cho ngay

\begin{equation}
  \norm{\jac{\vect{a}_t}{\vect{a}_k}} \le \eta^{\,t-k}.
  \label{eq:ch03-chan-mu}
\end{equation}

Số hạng dài hạn, tức những \(k\) cách \(t\) rất xa, tắt theo cấp số nhân. Chứng
minh dừng ở đó, sau bốn dòng.

Với tanh thì \(\gamma = 1\), nên ngưỡng là 1. Với sigmoid thì \(\gamma = 1/4\),
nên ngưỡng là 4: một mạng dùng sigmoid chịu được ma trận hồi quy lớn hơn nhiều
trước khi gradient bắt đầu tắt, đổi lại nó tắt bằng một cơ chế khác, vì
\(\sigma'\) của sigmoid không bao giờ vượt \(1/4\) ngay từ đầu.

\subsection*{Điều kiện cần để gradient nổ}

\index{exploding gradient!điều kiện cần}%
Lật ngược lập luận trên: nếu gradient nổ thì bất đẳng thức~\eqref{eq:ch03-chan-mu}
không được phép đúng, nên \(\norm{\mat{W}_{hh}}\) phải lớn hơn \(1/\gamma\).
Đây là điều kiện \emph{cần}, không phải \tn{điều kiện đủ}{sufficient condition}, và sự bất đối xứng ấy có
thật chứ không phải một câu rào đón. Vượt ngưỡng không có nghĩa là gradient sẽ
nổ.

Tôi đo thử. Bốn ma trận hồi quy chuẩn tắc với \tn{bán kính phổ}{spectral radius} đặt trước, 64 đơn vị
ẩn, tanh, chạy 100 bước, rồi lấy chuẩn phổ của tích
trong~\eqref{eq:ch03-tich-jacobian} theo từng khoảng cách:

\begin{minted}[fontsize=\footnotesize]{text}
radius  norm   l=1         l=10        l=50        l=100      rate
0.500   0.500  5.0000e-01  9.7532e-04  8.8705e-16  7.8786e-31 0.5000
0.900   0.900  9.0000e-01  3.4200e-01  5.0259e-03  2.5902e-05 0.9000
1.000   1.000  1.0000e+00  9.5614e-01  7.6328e-01  6.2852e-01 0.9961
1.200   1.200  1.2000e+00  4.9458e+00  2.4703e+01  8.2505e+01 1.0244
\end{minted}

\tn{Ma trận chuẩn tắc}{normal matrix} là trường hợp dễ: với chúng, bán kính phổ
và chuẩn phổ là cùng
một số, nên chặn~\eqref{eq:ch03-chan-tung-thua-so} chặt và lý thuyết mô tả đúng
thực tế. Hai dòng đầu khớp hoàn hảo: tốc độ đo được là 0.5000 và 0.9000, đúng
bằng bán kính phổ.

\begin{measured}{tốc độ tắt so với bán kính phổ}
Hai dòng dưới mới là chỗ đáng nhìn.

Ở bán kính phổ đúng bằng 1, tích vẫn tắt, chỉ là chậm: tốc độ đo được 0.9961
chứ không phải 1, và sau 100 bước chuẩn còn 0.62852. Lý do nằm ở
\(\sigma'\): tanh có \(\sigma'(x) = 1\) đúng một điểm duy nhất, tại gốc. Trạng
thái vừa rời khỏi gốc là \(\gamma = 1\) trở thành một chặn không bao giờ đạt
được.

Ở bán kính phổ 1.2, tích tăng, nhưng tốc độ là 1.0244 chứ không phải 1.2. Sau
100 bước nó đạt 82.505, trong khi \(1.2^{100}\) vào khoảng 8.2818 nhân mười mũ
bảy, tức lớn hơn số đo được chừng một triệu lần. Trạng thái lớn dần, tanh
\tn{bão hòa}{saturation}, \(\sigma'\) sụt về gần không, và
mức tăng tự chặn lại. Đây chính là lý do điều kiện nổ chỉ là \tn{điều kiện cần}{necessary condition}.

Đo bằng \code{experiments/ch03\_decay.py} tại tag \code{ch03}.
\end{measured}

Bão hòa tự chặn là một tin tốt và một tin xấu. Tin tốt: gradient nổ hiếm khi
nổ mãi. Tin xấu: cơ chế chặn nó cũng là cơ chế xóa trí nhớ, vì
\(\sigma'\) nhỏ ở mọi bước sau đó. Chương~\ref{ch:lstm} sẽ sửa đúng chỗ này.
